\textbf{Phần 1. Trắc nghiệm lựa chọn}
\setcounterpageref{ex}{0}

\begin{ex} %[3H9N53]
Chu vi hình vuông có cạnh $7\text{ cm}$ là
\choice
{$49\text{ cm}$}
{\True $28\text{ cm}$}
{$35\text{ cm}$}
{$21\text{ cm}$}
\loigiai{
$7 \times 4 = 28\text{ cm}$.\\
\textbf{Đáp án của câu hỏi là B}
}
\end{ex}

\begin{ex} %[3H9N53]
Diện tích hình vuông có cạnh $7\text{ cm}$ là
\choice
{$28\text{ cm}^2$}
{$35\text{ cm}^2$}
{$56\text{ cm}^2$}
{\True $49\text{ cm}^2$}
\loigiai{
$7 \times 7 = 49\text{ cm}^2$.\\
\textbf{Đáp án của câu hỏi là D}
}
\end{ex}

\begin{ex} %[3H9H53]
Hình chữ nhật có chiều dài $10\text{ cm}$, chiều rộng $4\text{ cm}$. Chu vi của nó là
\choice
{$40\text{ cm}$}
{\True $28\text{ cm}$}
{$14\text{ cm}$}
{$20\text{ cm}$}
\loigiai{
$(10 + 4) \times 2 = 28\text{ cm}$.\\
\textbf{Đáp án của câu hỏi là B}
}
\end{ex}

\begin{ex} %[3H9H53]
Hình chữ nhật có chiều dài $10\text{ cm}$, chiều rộng $4\text{ cm}$. Diện tích của nó là
\choice
{$28\text{ cm}^2$}
{$20\text{ cm}^2$}
{\True $40\text{ cm}^2$}
{$30\text{ cm}^2$}
\loigiai{
$10 \times 4 = 40\text{ cm}^2$.\\
\textbf{Đáp án của câu hỏi là C}
}
\end{ex}

\begin{ex} %[3H9H53]
Hình vuông có diện tích $25\text{ cm}^2$. Cạnh hình vuông đó là
\choice
{\True $5\text{ cm}$}
{$4\text{ cm}$}
{$6\text{ cm}$}
{$10\text{ cm}$}
\loigiai{
Vì $5 \times 5 = 25$ nên cạnh là $5\text{ cm}$.\\
\textbf{Đáp án của câu hỏi là A}
}
\end{ex}

\begin{ex} %[3H9H53]
Hình chữ nhật có diện tích $36\text{ cm}^2$, chiều rộng $4\text{ cm}$. Chiều dài của hình là
\choice
{$8\text{ cm}$}
{$10\text{ cm}$}
{\True $9\text{ cm}$}
{$12\text{ cm}$}
\loigiai{
$36 : 4 = 9\text{ cm}$.\\
\textbf{Đáp án của câu hỏi là C}
}
\end{ex}

\begin{ex} %[3H9V53]
Một hình vuông có chu vi $32\text{ cm}$. Diện tích hình vuông đó là
\choice
{$32\text{ cm}^2$}
{$48\text{ cm}^2$}
{$81\text{ cm}^2$}
{\True $64\text{ cm}^2$}
\loigiai{
Cạnh: $32 : 4 = 8\text{ cm}$.\\
Diện tích: $8 \times 8 = 64\text{ cm}^2$.\\
\textbf{Đáp án của câu hỏi là D}
}
\end{ex}

\begin{ex} %[3H9V53]
Một hình chữ nhật có chiều dài $14\text{ cm}$, chiều rộng bằng nửa chiều dài. Diện tích hình chữ nhật đó là
\choice
{\True $98\text{ cm}^2$}
{$84\text{ cm}^2$}
{$102\text{ cm}^2$}
{$70\text{ cm}^2$}
\loigiai{
Rộng: $14 : 2 = 7\text{ cm}$.\\
Diện tích: $14 \times 7 = 98\text{ cm}^2$.\\
\textbf{Đáp án của câu hỏi là A}
}
\end{ex}

\begin{ex} %[3H9V53]
Ghép 2 hình vuông cạnh $4\text{ cm}$ lại thành một hình chữ nhật. Diện tích hình chữ nhật đó là
\choice
{$24\text{ cm}^2$}
{$36\text{ cm}^2$}
{\True $32\text{ cm}^2$}
{$16\text{ cm}^2$}
\loigiai{
$4 \times 4 \times 2 = 32\text{ cm}^2$.\\
\textbf{Đáp án của câu hỏi là C}
}
\end{ex}

\begin{ex} %[3H9V53]
Tính diện tích hình chữ nhật có nửa chu vi là $16\text{ cm}$, chiều dài là $10\text{ cm}$:
\choice
{$50\text{ cm}^2$}
{$70\text{ cm}^2$}
{$80\text{ cm}^2$}
{\True $60\text{ cm}^2$}
\loigiai{
Rộng: $16 - 10 = 6\text{ cm}$.\\
Diện tích: $10 \times 6 = 60\text{ cm}^2$.\\
\textbf{Đáp án của câu hỏi là D}
}
\end{ex}

\begin{ex} %[3H9V53]
Một viên gạch men hình vuông cạnh $20\text{ cm}$. Diện tích của viên gạch đó là
\choice
{$80\text{ cm}^2$}
{\True $400\text{ cm}^2$}
{$200\text{ cm}^2$}
{$40\text{ cm}^2$}
\loigiai{
$20 \times 20 = 400\text{ cm}^2$.\\
\textbf{Đáp án của câu hỏi là B}
}
\end{ex}

\begin{ex} %[3H9C53]
Một mảnh đất hình chữ nhật có chu vi $48\text{ m}$, nếu tăng chiều rộng thêm $4\text{ m}$ thì mảnh đất trở thành hình vuông. Diện tích mảnh đất ban đầu là
\choice
{\True $140\text{ m}^2$}
{$120\text{ m}^2$}
{$144\text{ m}^2$}
{$128\text{ m}^2$}
\loigiai{
Nửa chu vi: $48 : 2 = 24\text{ m}$.\\
Chiều dài hơn chiều rộng $4\text{ m}$.\\
Chiều rộng: $(24 - 4) : 2 = 10\text{ m}$.\\
Chiều dài: $10 + 4 = 14\text{ m}$.\\
Diện tích: $14 \times 10 = 140\text{ m}^2$.\\
\textbf{Đáp án của câu hỏi là A}
}
\end{ex}

\textbf{Phần 2. Trắc nghiệm đúng sai}
\setcounterpageref{ex}{0}

\begin{ex} %[3H9H53]
Xét tính đúng sai về chu vi và diện tích:
\choiceTF
{\True Chu vi hình vuông cạnh $6\text{ cm}$ là $24\text{ cm}$.}
{\True Diện tích hình vuông cạnh $6\text{ cm}$ là $36\text{ cm}^2$.}
{Hình vuông có chu vi lớn hơn thì chắc chắn có diện tích bé hơn.}
{\True Cạnh hình vuông càng lớn thì cả chu vi và diện tích đều càng lớn.}
\loigiai{
\begin{itemize}
    \item Chu vi: $6 \times 4 = 24\text{ cm}$. Đúng.
    \item Diện tích: $6 \times 6 = 36\text{ cm}^2$. Đúng.
    \item Chu vi lớn hơn thì cạnh lớn hơn nên diện tích lớn hơn. Ý 3 Sai, Ý 4 Đúng.
\end{itemize}
\begin{itemize}
    \item a) \textbf{Đúng}.
    \item b) \textbf{Đúng}.
    \item c) \textbf{Sai}.
    \item d) \textbf{Đúng}.
\end{itemize}
}
\end{ex}

\begin{ex} %[3H9V53]
Một hình chữ nhật có chiều dài $9\text{ cm}$, chiều rộng $4\text{ cm}$. Xét tính đúng sai:
\choiceTF
{\True Chu vi hình chữ nhật là $26\text{ cm}$.}
{\True Diện tích hình chữ nhật là $36\text{ cm}^2$.}
{Diện tích hình chữ nhật này bằng diện tích hình vuông cạnh $5\text{ cm}$.}
{\True Diện tích hình chữ nhật này bằng diện tích hình vuông cạnh $6\text{ cm}$ (vì $6 \times 6 = 36$).}
\loigiai{
\begin{itemize}
    \item Chu vi: $(9 + 4) \times 2 = 26\text{ cm}$. Đúng.
    \item Diện tích: $9 \times 4 = 36\text{ cm}^2$. Đúng.
    \item Bằng hình vuông cạnh $6\text{ cm}$ ($36\text{ cm}^2 \neq 25\text{ cm}^2$). Ý 3 Sai, Ý 4 Đúng.
\end{itemize}
\begin{itemize}
    \item a) \textbf{Đúng}.
    \item b) \textbf{Đúng}.
    \item c) \textbf{Sai}.
    \item d) \textbf{Đúng}.
\end{itemize}
}
\end{ex}

\textbf{Phần 3. Trả lời ngắn}
\setcounterpageref{ex}{0}

\begin{ex} %[3H9H53]
Tính diện tích hình chữ nhật có chiều dài $11\text{ cm}$, chiều rộng $5\text{ cm}$.
\shortans{55}
\loigiai{
\textbf{Đáp số:} $55$\\
$11 \times 5 = 55\text{ cm}^2$.
}
\end{ex}

\begin{ex} %[3H9H53]
Tính diện tích hình vuông có cạnh $9\text{ cm}$.
\shortans{81}
\loigiai{
\textbf{Đáp số:} $81$\\
$9 \times 9 = 81\text{ cm}^2$.
}
\end{ex}

\begin{ex} %[3H9V53]
Một hình vuông có diện tích $49\text{ cm}^2$. Tính chu vi của hình vuông đó theo xăng-ti-mét.
\shortans{28}
\loigiai{
\textbf{Đáp số:} $28$\\
Cạnh là $7\text{ cm}$. Chu vi: $7 \times 4 = 28\text{ cm}$.
}
\end{ex}

\begin{ex} %[3H9V53]
Một hình chữ nhật có diện tích $72\text{ cm}^2$, chiều dài $9\text{ cm}$. Tính chiều rộng hình chữ nhật theo xăng-ti-mét.
\shortans{8}
\loigiai{
\textbf{Đáp số:} $8$\\
$72 : 9 = 8\text{ cm}$.
}
\end{ex}

\begin{ex} %[3H9V53]
Một hình chữ nhật có chu vi $30\text{ cm}$, chiều dài $9\text{ cm}$. Tính diện tích hình chữ nhật theo xăng-ti-mét vuông.
\shortans{54}
\loigiai{
\textbf{Đáp số:} $54$\\
Nửa chu vi: $15\text{ cm}$. Rộng: $15 - 9 = 6\text{ cm}$. Diện tích: $9 \times 6 = 54\text{ cm}^2$.
}
\end{ex}

\begin{ex} %[3H9C53]
Ghép 4 hình vuông nhỏ cạnh $3\text{ cm}$ thành một hình vuông lớn. Tính diện tích hình vuông lớn theo xăng-ti-mét vuông.
\shortans{36}
\loigiai{
\textbf{Đáp số:} $36$\\
$3 \times 3 \times 4 = 36\text{ cm}^2$ (hoặc cạnh lớn $6\text{ cm} \Rightarrow 6 \times 6 = 36$).
}
\end{ex}
